% 英文源分章 PDF 第 11—14 页；第 9.7—9.8 节。
\section{变量重要性度量的 CV-TMLE}

对第 9.3 节所述算法的修改包括两个小变化：在训练样本上估计 $Q_0$ 和 $g_0$（也就是上文用来确定 $a_L,a_H$ 的同一份样本）；对巧妙协变量前的系数 $\epsilon$ 的估计，不仅基于相应的估计样本，而是基于整个样本。因此，将式 (9.12) 修改为
\begin{equation}
\begin{aligned}
\psi^{CV-TMLE}_{n,v}
={}&\frac1{n_V}\sum_{i:Z_i=v}\bigl\{Q_{n,v^c,\epsilon_n}(a_H(P_{n,v^c}),W_i)\\
&\hspace{7.3em}-Q_{n,v^c,\epsilon_n}(a_L(P_{n,v^c}),W_i)\bigr\}.
\end{aligned}
\tag{9.19}
\end{equation}
这里，式 (9.12) 与式 (9.19) 的区别是：(1) $Q_{n,v}=\hat Q(P_{n,v})$ 和 $g_{n,v}=\hat g(P_{n,v})$ 分别改为 $\hat Q(P_{n,v^c})$ 和 $\hat g(P_{n,v^c})$；(2) $\epsilon$ 从 $\epsilon_{n,v}$（仅在估计样本 $P_{n,v}$ 上估计）改为 $\epsilon_n$（在整个样本上估计）。这也要求将巧妙协变量的定义 (9.11) 改为
\begin{equation}
H_{P_{n,v^c}}(A,W;g_{n,v^c})
=\frac{I(A=a_H(P_{n,v^c}))}{g_{n,v^c}(a_H(P_{n,v^c})\mid W)}
-\frac{I(A=a_L(P_{n,v^c}))}{g_{n,v^c}(a_L(P_{n,v^c})\mid W)}.
\tag{9.20}
\end{equation}
因此，$\epsilon_n$ 是按式 (9.10) 进行逻辑回归的结果：将 $Y_i$ 对协变量 $H_{P_{n,v^c}}(A_i,W_i;g_{n,v^c})$ 回归，以 $\operatorname{Logit}Q_{n,v^c}(A_i,W_i)$ 为偏移项（offset），$i=1,\ldots,n$。这个 $\epsilon_n$ 为每个 $v=1,\ldots,V$ 给出更新后的 $Q_{n,v^c,\epsilon_n}$。最后，目标数据自适应变量重要性度量 $\psi_{0,n}$ 的估计量为
\begin{equation}
\hat\Psi(P_n)=\frac1V\sum_{v=1}^V\psi^{CV-TMLE}_{n,v}.
\tag{9.21}
\end{equation}
另外，对每个 $v=1,\ldots,V$，将影响曲线估计为
\[
\begin{aligned}
D^*_{P_{n,v^c}}(O)
={}&\left[\frac{I(A=a_{H,n,v^c})}{g_{n,v^c}(a_{H,n,v^c}\mid W)}
-\frac{I(A=a_{L,n,v^c})}{g_{n,v^c}(a_{L,n,v^c}\mid W)}\right]
\{Y-Q_{n,v^c,\epsilon_n}(A,W)\}\\
&+\bigl[Q_{n,v^c,\epsilon_n}(a_{H,n,v^c},W)
-Q_{n,v^c,\epsilon_n}(a_{L,n,v^c},W)\bigr]
-\psi^{CV-TMLE}_{n,v}.
\end{aligned}
\]
最后，作出这些修改后，可以按照式 (9.15) 中的同样方法进行推断。

这为原来同一个平均数据自适应参数提供了另一种估计量，但根据第 A.1 节的定理，在更大的模型中也能得到便于应用的渐近性质。从直觉上说，与第 9.3 节介绍的原算法相比，估计所需的各个构成参数 $(Q_0,g_0,\epsilon)$ 如今都基于更大比例的数据来估计，因此应当能够获得稳健得多的估计。这一 CV-TMLE 是一个复杂算法，但幸运的是，针对这里的应用已有相应的 R 软件包，下面将进一步介绍。

\section{数据自适应 VIM 的软件：varImpact}

我们提供了一个名为 \texttt{varImpact} 的软件包，用 R 编程语言实现（R Development Core Team，2016）。\texttt{varImpact} 可从综合 R 档案网络（Comprehensive R Archive Network，CRAN）和 Github 获取。\footnote{\url{http://github.com/ck37/varImpact/}。}\texttt{varImpact} 实现了本章介绍的变量重要性算法，还提供数据清理、报告生成及其他相关功能，以便在实际数据集中进行变量重要性分析。下面介绍 \texttt{varImpact} 算法的每一步及其并行化方法。

\noindent\textbf{1. 预处理。}\texttt{varImpact} 首先对数据集进行预处理，分别处理因子变量和数值变量。如果某个变量的缺失比例超过阈值（默认为 50\%），或者变异不足，则将其从分析中删除。这样可以防止过拟合，并将分析集中于测量覆盖率合理的变量。检查变异的步骤会分析每个变量的密度，删除第 10 百分位数与第 90 百分位数相等的变量。当数值变量作为关注变量（而不是调整变量）时，只要其不同取值超过十个，就将其离散化为十个分位数组。为其余调整变量生成缺失指示变量，将这些指示变量纳入调整集，并用 $k$ 近邻法填补缺失值；软件也支持中位数插补和零值替代插补。将数据集划分为 CV-TMLE 的 $V$ 折，建议取 $V=10$，以充分利用 CV-TMLE 的功效。当结局为二元变量时，按结局分层进行划分，以使功效最大化。对所分析的每个变量使用同一划分。

\noindent\textbf{2. 对每个变量构造观察性研究。}现在，可以为每个变量构造一项数据自适应的观察性研究，在控制其余所有调整变量（协变量和缺失指示变量）的同时，估计该变量最具影响的变化如何影响结局。依次将每个变量看作多值处理或干预。首先，使用训练数据估计每个处理水平下的潜在结局均值。然后，使用留出的测试数据，将估计结局均值最高与最低的两个水平之间的均值差，作为变量重要性度量的估计。这就是将当前变量设为“最佳”水平相对于“最差”水平时的平均处理效应估计。在分析每个变量时，\texttt{varImpact} 遍历 CV-TMLE 的各折，将每一折视为测试集，其余各折视为训练集。

\noindent\textbf{3. 逐折分析。}每一次折迭代（CV-TMLE 至少需要两折）的分析包含若干步骤。通过带惩罚的直方图密度估计（Rozenholc 等，2010；Mildenberger 等，2009），进一步合并离散化数值变量的区间，以避免单元格样本量过小。软件还提供一个选项：对调整变量进行层次聚类，然后选出最具代表性的前十个变量（中心点，medoids），所用聚类算法为 HOPACH（van der Laan 和 Pollard，2003）。这种降维能够大幅加快计算，但也容易导致偏差或功效损失。对因子变量，我们还检查其因子水平与二元结局变量交叉列联时，每个单元格的最小观测数。可以跳过单元格样本量过小的协变量，以节省计算时间并减轻过拟合。

然后，在训练集中，用 TMLE 估计每个水平或区间 $a\in\mathcal A$ 的调整后结局均值。\footnote{与 TMLE R 软件包（Gruber 和 van der Laan，2012a）一样，将估计的 $g$ 截断到 $[0.025,0.975]$ 内。同样按照 TMLE R 软件包的做法，我们对 $Q$ 和 $g$ 都采用非负最小二乘作为元学习器。}特定于 $a$ 的调整后结局均值记为 $\theta_0(a)\equiv E_{0,W}\{Q_0(a,W)\}$，特定训练样本 $v^c$ 上的估计记为 $\hat\theta_{P_{n,v^c}}(a)$。我们找出分别与最高和最低结局均值对应的区间／水平：
\[
\begin{aligned}
a_L(P_{n,v^c})&=\argmin_{a\in\mathcal A}\hat\theta_{P_{n,v^c}}(a),\\
a_H(P_{n,v^c})&=\argmax_{a\in\mathcal A}\hat\theta_{P_{n,v^c}}(a).
\end{aligned}
\]
该变量取值处于“高”区间 $(a_H)$ 的观测实际上构成处理组，处于“低”区间 $(a_L)$ 的观测构成对照组。保存训练集上相应的结局回归 $(Q)$ 和倾向评分 $(g)$ 的 \texttt{SuperLearner} 模型拟合。随后，在相应测试集上应用保存的 $Q$ 和 $g$ 模型拟合，对估计样本作出所需预测。

\noindent\textbf{4. 汇集各折结果。}在每个 $V$ 折训练样本上估计出所有干扰参数以及 $a_L,a_H$ 后，就可以基于完整数据集（即 $V$ 个测试样本的并集）估计 $\epsilon$。实际操作中，我们分别为 $a_L,a_H$ 两个水平构造各自的巧妙协变量（式 (9.20)），并通过逻辑回归估计二维扰动系数 $\epsilon_n$。随后对预测结局进行扰动，靶向这两个水平下的结局均值，再对每个测试集分别计算该次划分特定的靶向结局均值及相关影响曲线值。每一折内，高水平 $(a_H)$ 的靶向均值与低水平 $(a_L)$ 的靶向均值之差，就是该折特定的平均处理效应（average treatment effect，ATE；式 (9.19)）。类似地，每一折内，$a_H$ 的影响曲线与 $a_L$ 的影响曲线之差，给出 ATE 的影响曲线。特定于该折的 ATE 的样本方差估计 $\sigma_{n,v}^2$，就是特定于该折的影响曲线的样本方差估计 $\sigma_{n,v}^2$（式 (9.18)）。

\noindent\textbf{5. 点估计与推断。}为生成最终结果，我们合并各折的参数估计。将各折 ATE 的均值作为 ATE 的点估计（式 (9.21)）：$\psi_n=\operatorname{Ave}\{\hat\psi_{P_{n,v^c}}(P_{n,v})\}$；将各折样本方差的均值作为样本方差的估计：$\sigma_n^2=\operatorname{Ave}\{\sigma_{n,v}^2\}$。估计标准误为 $\sqrt{\sigma_n^2/n}$。我们报告基于正态分布的置信区间，以及针对原假设 $H_0:\psi_{n,0}\leq0$ 的单侧 $p$ 值。双侧检验在这里并不适当，因为处理水平 $a_H$ 和 $a_L$ 正是为了得到正的处理效应而选择的。任何负的处理效应估计，都表明该程序利用所识别的水平未能为这一特定变量发现处理效应。

\noindent\textbf{6. 报告结果。}在最后的报告阶段，我们进行多重比较调整，确定最终的统计显著性，并标记在 $V$ 折间结果不一致的变量。我们的多重比较调整通过 Benjamini–Hochberg 程序（Benjamini 和 Hochberg，1995）控制错误发现率（false discovery rate，FDR）。FDR 调整后的 $p$ 值小于 0.05 的变量被判定为具有统计显著性。对某个分类变量，若在 $V$ 折中均选择了相同的 $a_L$ 和 $a_H$ 水平，就称该变量的结果具有一致性。分析者可以放宽这一标准，例如仅要求至少一定比例的折选择相同的 $a_L$ 和 $a_H$ 水平。对于数值变量，一致性定义为所有 $V$ 折在低、高分位数的方向上相同。换言之，一致的结果是每个 CV-TMLE 折的 $a_L$ 都是低于 $a_H$ 的分位数，或者每个 CV-TMLE 折的 $a_L$ 都是高于 $a_H$ 的分位数。变量按 $p$ 值升序排序，列出其排名、参数估计、未经调整的 $p$ 值、FDR 调整后的 $p$ 值以及 95\% 置信区间。

\noindent\textbf{7. 并行化。}对数据集中的每个变量开展半参数观察性研究，计算需求很大。当采用为准确估计结局和倾向评分所必需的复杂 \texttt{SuperLearner} 候选库时，情况尤为如此。CV-TMLE 进一步增加所需计算，因为它实际上对每个变量重复开展多次观察性研究。为解决这一问题，\texttt{varImpact} 通过 \texttt{future} 软件包支持并行化，可以无缝使用一台机器的多个核心，或者集群中的多台机器。这能够大幅缩短分析所需的总历时，即“墙钟”时间。
